% -----------------------------------------------------------------------------
% Dirac/Clifford identities, proved after gamma matrices are constructed
% -----------------------------------------------------------------------------

Poincar\'e 表示与 Clifford 代数确定 $\gamma^\mu$ 以后，后续散射与圈图所需的 Dirac 矩阵代数不再是独立公式表，而都来自
$\{\gamma^\mu,\gamma^\nu\}=2\eta^{\mu\nu}\id(4)$、矩阵迹的循环性以及 $\gamma^5$ 的定义。本段只在正文列出会被反复引用的目标恒等式；逐次反对易和指标缩并统一放在本结构末尾的推导区。

最简单的 Lorentz 二阶张量迹只能正比于 $\eta^{\mu\nu}$。取 $\mu=\nu=0$ 并使用 $(\gamma^0)^2=\id(4)$ 固定系数后，得到
\begin{equation}
 \boxed{\Tr(\gamma^\mu\gamma^\nu)=4\eta^{\mu\nu}.}
 \label{eq:trace-two-gamma-r9}
\end{equation}
四个伽马矩阵的迹可把最左因子依次反对易到最右端，再利用\eqnrefs{eq:trace-two-gamma-r9} 关闭递推，结果为
\begin{equation}
 \boxed{
 \Tr(\gamma^\mu\gamma^\nu\gamma^\rho\gamma^\sigma)
 =4(\eta^{\mu\nu}\eta^{\rho\sigma}
 -\eta^{\mu\rho}\eta^{\nu\sigma}
 +\eta^{\mu\sigma}\eta^{\nu\rho}).}
 \label{eq:trace-four-gamma-r9}
\end{equation}
这两条迹恒等式是未偏振费米子散射把自旋求和转成 Lorentz 标量的基本工具。


当费米子链继续增长时，不需要为六个、八个乃至更多 $\gamma$ 矩阵分别记忆新公式。把最左端矩阵逐次反对易到最右端，再用迹的循环性关闭链条，可得到统一递推
\begin{equation}
 \boxed{
 T_{2n}^{\mu_1\cdots\mu_{2n}}
 =\sum_{j=2}^{2n}(-1)^j\eta^{\mu_1\mu_j}
 \Tr\!\left(\gamma^{\mu_2}\cdots\widehat{\gamma^{\mu_j}}\cdots\gamma^{\mu_{2n}}\right),}
 \qquad
 T_{2n}^{\mu_1\cdots\mu_{2n}}
 \equiv\Tr(\gamma^{\mu_1}\cdots\gamma^{\mu_{2n}}).
 \label{eq:trace-even-recursion-r68}
\end{equation}
这里帽号表示删除相应因子。另一方面，在严格四维且迹中不含额外 $\gamma^5$ 时，$\gamma^5$ 与每个普通 $\gamma^\mu$ 反对易，因此任意奇数长度的普通 Dirac 链满足
\begin{equation}
 \boxed{\Tr(\gamma^{\mu_1}\cdots\gamma^{\mu_{2n+1}})=0.}
 \label{eq:trace-odd-zero-r68}
\end{equation}
这两条式子分别负责把长偶数链递归缩短，以及在自旋求和中立即消去奇数链。涉及维数正规化中的真正手征迹时，$\gamma^5$ 的处理需按所采用的正规化方案另行定义。

对任意四矢量 $a^\mu,b^\mu,c^\mu$，定义斜线记号 $\slashed a\equiv\gamma^\mu a_\mu$。把外侧的 $\gamma^\mu$ 逐次穿过内部矩阵串，可得到长度为一、二、三的三条常用夹乘恒等式：
\begin{align}
 \boxed{\gamma^\mu\slashed a\gamma_\mu=-2\slashed a,}
 \label{eq:gamma-one-sandwich-r10}\\[2pt]
 \boxed{\gamma^\mu\slashed a\slashed b\gamma_\mu
 =4a\!\cdot b\,\id(4),}
 \label{eq:gamma-two-sandwich-r10}\\[2pt]
 \boxed{\gamma^\mu\slashed a\slashed b\slashed c\gamma_\mu
 =-2\slashed c\slashed b\slashed a.}
 \label{eq:gamma-three-sandwich-r10}
\end{align}
这些式子的输入只有四维 Clifford 代数；其中的数值系数 $-2,4,-2$ 会在维数正规化推广到 $d$ 维时发生相应改变，所以圈图计算若尚未取 $d\to4$，不能无条件使用四维版本。


在维数正规化中应保留 $d$ 维 Clifford 代数本身。由 $g^\mu{}_{\mu}=d$ 逐项收缩可得最常用的一阶夹乘
\begin{equation}
 \boxed{\gamma^\mu\gamma^\alpha\gamma_\mu=(2-d)\gamma^\alpha,}
 \label{eq:gamma-d-sandwich-r68}
\end{equation}
它在 $d=4$ 时才退化为 $-2\gamma^\alpha$。因此圈积分分子中的 $O(\epsilon)$ 代数不能在积分完成前被四维恒等式提前抹掉。

严格四维中，含一个 $\gamma^5$ 的最低阶非零迹由当前取向约定唯一固定为
\begin{equation}
 \boxed{\Tr(\gamma^5\gamma^\mu\gamma^\nu\gamma^\rho\gamma^\sigma)
 =-4\ii\epsilon^{\mu\nu\rho\sigma}.}
 \label{eq:gamma5-four-trace-r68}
\end{equation}
它把手征 Dirac 代数与 Levi--Civita 张量连接起来，后续轴矢量流与异常计算都从这一约定出发。

最后，复 $4\times4$ 矩阵空间有 $16$ 个复维度。由
$\{\id(4),\gamma^\mu,\sigma^{\mu\nu},\gamma^\mu\gamma^5,\gamma^5\}$
构造一组迹正交基 $\Gamma_A$，并以 $\Gamma^A$ 表示其对偶基，则矩阵指标的完备关系为
\begin{equation}
 \boxed{
 \frac14\sum_A
 (\Gamma_A)_{ij}(\Gamma^A)_{kl}
 =\delta_{il}\delta_{kj}.}
 \label{eq:dirac-matrix-completeness-r10}
\end{equation}
Fierz 重排只是把这条 $16$ 维矩阵完备关系插入两组旋量指标之间，再改变指标配对；因此 Fierz 重排完全由有限维 Dirac 矩阵空间的完备性与指标重配得到。


例如，对普通复数值外线旋量，标量双线性的一种常用重排为
\begin{align}
 (\bar u_1u_2)(\bar u_3u_4)
 =\frac14\Big[&
 (\bar u_1u_4)(\bar u_3u_2)
 +(\bar u_1\gamma^\mu u_4)(\bar u_3\gamma_\mu u_2)
 \nonumber\\
 &+\frac12(\bar u_1\sigma^{\mu\nu}u_4)(\bar u_3\sigma_{\mu\nu}u_2)
 -(\bar u_1\gamma^\mu\gamma^5u_4)(\bar u_3\gamma_\mu\gamma^5u_2)
 \nonumber\\
 &+(\bar u_1\gamma^5u_4)(\bar u_3\gamma^5u_2)\Big].
 \label{eq:fierz-scalar-r68}
\end{align}
若 $u_i$ 换成 Grassmann 场，还必须另外计入为改变场次序而产生的统计负号；矩阵完备系数与费米统计符号是两件不同的事。

同一组迹恒等式还给出 Lorentz 生成元矩阵的两个常用收缩：
\begin{equation}
 \boxed{\Tr(\sigma^{\mu\nu}\sigma^{\rho\sigma})
 =4(\eta^{\mu\rho}\eta^{\nu\sigma}-\eta^{\mu\sigma}\eta^{\nu\rho}),
 \qquad
 \sigma^{\mu\nu}\sigma_{\mu\nu}=12\id(4).}
 \label{eq:sigma-sigma-trace-r68}
\end{equation}
在当前 $\eta=\mathrm{diag}(1,-1,-1,-1)$、$\epsilon^{0123}=+1$ 约定下，$\gamma^5$ 把反对称张量映射到其 Hodge 对偶：
\begin{equation}
 \boxed{\sigma^{\mu\nu}\gamma^5
 =\frac{\ii}{2}\epsilon^{\mu\nu\rho\sigma}\sigma_{\rho\sigma}.}
 \label{eq:sigma-gamma5-dual-r68}
\end{equation}
这些结果分别用于四费米子算符换基、Pauli 形状因子以及手征张量结构。
